\section{Contoh Program}
\label{sec:contoh-bab5}

Bagian ini menyajikan tiga contoh. Dua contoh pertama dapat langsung
dirakit dan dijalankan, sedangkan contoh ketiga adalah potongan
singkat untuk memperlihatkan perilaku register pada tingkat byte.

\subsection{Contoh 1: Memprediksi Flag Sebelum Menjalankan Program}

Program pada Kode~\ref{lst:prediksi-flag} menjalankan dua
penjumlahan yang sengaja dipilih untuk memicu kondisi \eng{flag}
berbeda. Sebelum program dijalankan, mahasiswa diminta memprediksi
nilai \flag{CF}, \flag{OF}, \flag{SF}, dan \flag{ZF} pada
Tabel~\ref{tab:prediksi-flag}, lalu memeriksa hasil prediksinya
melalui \eng{Single Step}.

\lstinputlisting[
    language=MASM,
    caption={Program \file{prediksi-flag.asm} untuk mengamati
    perubahan \eng{flag}.},
    label={lst:prediksi-flag}
]{\CodeDir/prediksi-flag.asm}

\begin{table}[H]
    \centering
    \caption{Lembar prediksi \eng{flag} untuk
    Kode~\ref{lst:prediksi-flag}. Isi kolom terakhir berdasarkan
    pengamatan pada emulator.}
    \label{tab:prediksi-flag}
    \small
    \begin{tabular}{@{}clccccc@{}}
        \toprule
        \textbf{Kasus} & \textbf{Operasi} & \flag{CF} & \flag{OF} &
        \flag{SF} & \flag{ZF} & \textbf{Hasil (\reg{AX})} \\
        \midrule
        1 & \hex{007F} $+$ \hex{0001} & \dots & \dots & \dots &
        \dots & \dots \\
        \addlinespace
        2 & \hex{8000} $+$ \hex{8000} & \dots & \dots & \dots &
        \dots & \dots \\
        \bottomrule
    \end{tabular}
\end{table}

Penjelasan potongan penting pada program tersebut adalah sebagai
berikut.

\begin{itemize}
    \item \lstasm|MOV AX, 007Fh| memuat $+127$ ke akumulator tanpa
          mengubah \eng{flag}.
    \item \lstasm|ADD AX, BX| adalah instruksi yang mengubah
          \eng{flag}. Kedua operand bernilai positif, tetapi
          hasilnya \hex{0080} yang bit tandanya menyala. Inilah
          \eng{overflow} bertanda, sehingga \flag{OF} bernilai 1.
    \item \lstasm|MOV hasil1, AX| menyimpan salinan hasil ke segmen
          data. Karena \asm{MOV} tidak menyentuh \eng{flag}, nilai
          \eng{flag} hasil penjumlahan masih dapat diamati pada
          \eng{debugger} setelah baris ini dijalankan.
    \item Kasus kedua memakai nilai \hex{8000} pada kedua operand.
          Hasilnya nol dengan simpanan keluar, sehingga
          \flag{CF} dan \flag{ZF} bernilai 1 sekaligus.
\end{itemize}

\begin{session}[caption={Keluaran program \file{prediksi-flag.asm}.},
label={lst:keluaran-prediksi-flag}]
Pengamatan flag selesai.
\end{session}

\subsection{Contoh 2: Memilih Pesan Berdasarkan Flag}

Contoh pertama memperlihatkan bagaimana \eng{flag} terbentuk.
Contoh kedua, pada Kode~\ref{lst:pesan-flag}, memperlihatkan
bagaimana \eng{flag} \emph{dipakai}. Program menjumlahkan
\hex{F0} dan \hex{10} pada register \reg{AL} lalu mencetak pesan yang
berbeda sesuai keadaan \eng{flag}.

\lstinputlisting[
    language=MASM,
    caption={Program \file{pesan-flag.asm} memilih pesan berdasarkan
    \eng{flag} hasil \lstasm{ADD AL, BL}.},
    label={lst:pesan-flag}
]{\CodeDir/pesan-flag.asm}

Bagian yang perlu diperhatikan adalah urutan pemeriksaan
\eng{flag}. Setiap kali \lstasm|ADD AL, BL| selesai, register
\reg{FLAGS} memuat seluruh penanda hasil operasi tersebut sekaligus.
Program kemudian memeriksanya satu per satu:

\begin{enumerate}
    \item \lstasm|JNC tanpa_carry| melompat bila \flag{CF} bernilai
          0. Bila \flag{CF} bernilai 1, instruksi berikutnya
          mencetak pesan carry.
    \item \lstasm|JNZ bukan_nol| melompat bila \flag{ZF} bernilai 0.
          Bila hasilnya nol, program mencetak pesan hasil nol.
    \item \lstasm|JNS tanpa_negatif| melompat bila \flag{SF} bernilai
          0. Bila bit tanda menyala, program mencetak pesan hasil
          negatif.
    \item \lstasm|CMP AL, 0| dipakai untuk membangkitkan kembali
          \eng{flag} berdasarkan isi \reg{AL} saat itu. Instruksi
          ini tidak menyimpan hasil pengurangan, hanya meninggalkan
          jejak pada \eng{flag}.
\end{enumerate}

Karena \hex{F0} ditambah \hex{10} menghasilkan \hex{100}, register
\reg{AL} menyimpan \hex{00} dengan \flag{CF} bernilai 1 dan
\flag{ZF} bernilai 1. Keluaran program tampak pada
Sesi~\ref{lst:keluaran-pesan-flag}.

\begin{session}[caption={Keluaran program \file{pesan-flag.asm}.},
label={lst:keluaran-pesan-flag}]
Terjadi carry, CF = 1.
Hasil penjumlahan nol, ZF = 1.
\end{session}

\begin{infobox}
Layanan \code{INT 21h} tidak merusak \eng{flag} hasil penjumlahan.
Instruksi \asm{INT} menyimpan \reg{FLAGS} ke \eng{stack} dan
\asm{IRET} mengembalikannya, sehingga pemeriksaan \flag{CF},
\flag{ZF}, dan \flag{SF} secara berurutan tetap sah.
\end{infobox}

\subsection{Contoh 3: Akses Byte Rendah dan Byte Tinggi}

Potongan pada Kode~\ref{lst:akses-byte} memperlihatkan bahwa satu
register 16 bit dapat diperlakukan sebagai dua register 8 bit yang
mandiri.

\begin{lstlisting}[
    language=MASM,
    caption={Mengakses byte rendah dan byte tinggi register
    \reg{AX}.},
    label={lst:akses-byte}
]
.MODEL SMALL
.STACK 100h

.DATA
    dummy DB 0          ; variabel penanda agar segmen data ada

.CODE
START:
    MOV AX, @DATA
    MOV DS, AX

    MOV AX, 1234h       ; AH = 12h, AL = 34h
    MOV BL, AH          ; BL = 12h (salin byte tinggi)
    MOV BH, AL          ; BH = 34h (salin byte rendah)

    ADD AL, BL          ; AL = 34h + 12h = 46h, AH tidak berubah
    MOV CL, AL          ; CL = 46h

    MOV AH, 4Ch
    MOV AL, 00h
    INT 21h

END START
\end{lstlisting}

Setelah seluruh instruksi di atas dijalankan, isi register adalah
\reg{AX} = \hex{1246}, \reg{BX} = \hex{3412}, dan \reg{CL} =
\hex{46}. Perhatikan bahwa \lstasm|ADD AL, BL| hanya mengubah
\reg{AL}; nilai \reg{AH} tetap \hex{12} karena \eng{flag} dan
\reg{AH} tidak tersentuh oleh operasi 8 bit pada \reg{AL}.
